\section{Limitations and Non-Capabilities}
\label{sec:limitations}

\subsection{Scope of the implementation}

Sirus is a selective protector, not a whole-program virtualizer. It accepts a
subset of x86-64, a subset of the ABI, and a subset of the platform surface. The
boundaries are deliberate: the decoder fails closed rather than guessing, and the
runtime fails closed rather than accessing host state.

\begin{table}[htbp]
\centering
\small
\caption{Principal limitations.}
\label{tab:limitations}
\begin{tabularx}{\textwidth}{@{}p{0.20\textwidth}X@{}}
\toprule
Area & Limitation \\
\midrule
Instruction set & Common integer operations, private-memory widths, the documented SSE2 and 128-bit VEX integer forms, \code{rep movs/stos} on private memory, and \code{bswap}. No AVX-256/AVX2, x87, MMX, floating point, or EVEX. \\
Control flow & Direct calls between selected functions; no external, indirect, or unresolved calls; no indirect jumps or jump tables; no cross-block flag liveness that cannot be proven. \\
Memory & Private 4 KiB stack and 2 KiB read-only data; no host pointers; no writable RIP-relative globals; no absolute host memory. \\
ABI & System V 6 register + 6 stack integer arguments; Win64 4 register + 6 stack integer arguments; one integer result; no floating-point arguments or results. \\
Function size & 16 KiB per function; 8 functions per program; 128 KiB payload; 256 blocks; 8192 IR operations. \\
Platform & Linux ELF64 x86-64 is the primary path. Windows PE32+ AMD64 is a selective leaf-function line. No Mach-O, ARM64, or 32-bit x86. \\
Security & No kernel anti-debug, DBI detection, trusted signing, formal Spectre proof, or proof of obfuscation security. The normal profile leaks timing. \\
Tooling & No IDA Pro, Binary Ninja, Ferrocene, or Triton in the local gate; no independent audit. \\
\bottomrule
\end{tabularx}
\end{table}

\subsection{Instruction-set boundary}

The supported subset is listed in Appendix~\ref{app:x86}. A function that calls
\code{libcore}, formats text, allocates, touches \code{std}, uses floating point,
or invokes an external library is generally not directly virtualizable. The
intended target is compact leaf logic: a license predicate, a checksum, a key
derivation step, or a proprietary arithmetic routine.

The decoder rejects a form rather than approximating it. This is why a function
that compiles and runs natively can still fail protection: the subset is smaller
than the native ISA. The user must inspect the error and either refactor the
function or choose a different target.

\subsection{ABI boundary}

The VM returns zero on any fail-closed condition. The host must consume the
protected value as data so that a wrong value causes a downstream failure. A
boolean-only interface that treats zero as "denied" is a poor fit: an attacker
can patch the caller or force the return. The recommended design derives a key,
checksum, or protocol field from the protected value and uses that result in a
subsequent operation.

The Windows line supports the Win64 integer ABI and preserves the callee-saved
XMM range through the shim. It does not support floating-point arguments,
language exception propagation, arbitrary imports, or kernel callbacks.

\subsection{Platform boundary}

The Linux ELF path is the primary production path. It supports ELF64 x86-64
executables and shared objects, with symbol, substring, range, or marker
selection. Fully stripped binaries require marker or range mode. The output
retains the host symbol table; the user can strip it separately if desired.

The Windows PE path is a selective leaf-function line. It injects one
read-execute section, patches RVA-relative entries, merges linked unwind
metadata, and passes Wine, native Windows, and MSVC \code{/guard:cf} fixtures.
It does not claim full Win64 virtualization, signing, imports/TLS edge cases, or
runtime exception propagation.

\subsection{Security boundary}

The safe tier does not implement kernel anti-debug, DBI detection, debug-register
inspection, RDTSC timing checks, or a native text checksum. Those controls require
unsafe OS interfaces and are confined to the optional hardened tier. The
hardened tier itself does not claim kernel-level or hypervisor-level detection.

The payload construction is an obfuscation and tamper-evidence layer, not a
standard authenticated cipher. The body stream is a keyed ARX stream; SipHash is
used for the manifest, integrity, and key derivation. The integrity value is
64-bit and the key shares are present in the generated code. An attacker with
full local control can rewrite the generated code and the payload together.

The normal profile is not constant-time. The explicit mode is a fixed-pass
oblivious subset, not a general constant-time compiler, and it does not model
speculative execution.

\subsection{Verification boundary}

Kani verifies bounded memory and arithmetic harnesses. Miri checks the VM library
tests. libFuzzer explores payloads, lifter inputs, and binary parsers. Ghidra,
radare2, angr, KLEE, GAMBA, SiMBA, MBA-Blast, dudect, and Cachegrind are run with
fixed budgets. None of these is a whole-product proof. The decoder is too large
for a useful CBMC run on this host. The symbolic probes time out rather than
prove unsatisfiability. The MBA study shows reductions, not irreducibility.

\subsection{What ``production prototype'' means}

The Linux pipeline is real and executed end-to-end. The safety gates, fuzzing,
Miri, Kani, differential builds, tamper and bypass tests, external analysis, and
Windows CI pass. The remaining gaps are explicit research, platform, and tooling
gaps, not hidden demo shortcuts. A release to untrusted third parties should
require an independent audit, a formal review of the custom payload
construction, and platform-specific validation beyond the tested fixtures.
